\subsubsection*{命题 3.5.}
\underline{假设态射 $\varphi:\mathbb P(N)\to\check{\mathbb P}^{r}$ 一般非分歧(即不处处分歧). 则 $\varphi$ 诱导从 $\mathbb P(N)$ 到 $\check X$ 上的双有理态射. 更准确地说, $\check X$ 的以下三个开子集相同:}
\begin{enumerate}
\item[(a)] \underline{$\check X$ 光滑的最大开子集;}
\item[(b)] \underline{使 $\varphi$ 诱导同构 $\varphi^{-1}(V)\xrightarrow{\sim}V$ 的最大开子集 $V\subset\check X$;}
\item[(c)] \underline{由满足 $X\cap H$ 恰有一个奇点且该点为非退化二次奇点的超平面 $H$ 构成的开子集(参见 3.2).}
\end{enumerate}

\subsubsection*{推论 3.5.0.}
\underline{假设 $\operatorname{char}(k)\ne2$ 或 $\dim X$ 为偶数. 给定嵌入为 Lefschetz 嵌入, 当且仅当以下两个条件之一成立:}
\begin{enumerate}
\item[(a)] $\varphi$ \underline{一般非分歧;}
\item[(b)] $\dim\check X\leq r-2$.
\end{enumerate}

事实上, 充分性由判别条件 (3.2.4) 得出. 在情形 (a) 中, 命题 3.5 将 $F$ 等同于 $\check X$ 的奇点集, 从而
\[
\dim F<\dim\check X=r-1;
\]
在情形 (b) 中, 此判别条件显然成立. 反之, 假设 $D$ 为 Lefschetz 线束, 且 $\dim\check X=r-1$. 存在点 $H\in D\cap\check X$, 由 3.3 及定义 2.2(c) 可知 $\varphi$ 在 $\varphi^{-1}(H)$ 的各点处非分歧. 因而它一般非分歧.

\noindent\underline{命题 3.5 的证明.}
假设蕴含
\[
\dim\check X=\dim\mathbb P(N)=r-1,
\]
于是 3.2 表明 $\check X$ 为 $\check{\mathbb P}^{r}$ 中的不可约超曲面.

设 $R$ 为 $\varphi$ 的分歧轨迹. 由假设, 它是 $\mathbb P(N)$ 中维数至多 $r-2$ 的闭子集; 故 $\varphi(R)$ 为 $\check X$ 中维数至多 $r-2$ 的闭子集. 设 $V'$ 为 $\check X$ 的稠密开子集
\begin{equation}\tag{3.5.1}
V'=(\check X-\varphi(R))\cap
   (\check X-\operatorname{Sing}(\check X))
\end{equation}
并以 $F$ 表示其补集:
\begin{equation}\tag{3.5.1.1}
F=\varphi(R)\cup\operatorname{Sing}(\check X).
\end{equation}

首先通过构造逆态射 $\psi:V'\to\varphi^{-1}(V')$ 来证明 $\varphi$ 诱导同构 $\varphi^{-1}(V')\simeq V'$. 由 (3.1.1) 将 $\varphi^{-1}(V')$ 视作嵌入 $X\times V'$ 中, 从而 $\varphi$ 为第二投影. 于是构造 $\psi$ 等价于构造一个使图
\begin{equation}\tag{3.5.2}
\begin{tikzcd}[column sep=small]
\mathbb P(N) \arrow[dd]
  & {}\supset\varphi^{-1}(V') \arrow[r,"\varphi"]
  & V' \arrow[dd] & \\
  & & & \psi_1 \\
X \arrow[rr,hook] & & \mathbb P^r &
\end{tikzcd}
\end{equation}
交换的态射 $\psi_1:V'\to\mathbb P^r$; 然后置 $\psi(v)=(\psi_1(v),v)\in\mathbb P^r\times V'$.

对射影空间中的任意超曲面, Gauss 态射 (3.4.1) 在每个光滑点处都有定义. 将此应用于超曲面 $\check X\subset\check{\mathbb P}^{r}$. 由于 $V'$ 位于 $\check X$ 的光滑轨迹 $V$ 内, Gauss 态射给出
\begin{equation}\tag{3.5.3}
\psi_1:V'\longrightarrow
\check{\check{\mathbb P}}^{r}\simeq\mathbb P^r.
\end{equation}
尚需检验 (3.5.2) 交换. 设 $(x_0,H_0)$ 为 $\varphi^{-1}(V')$ 的闭点. 须证明将 $x_0$ 视作 $\check{\mathbb P}^{r}$ 中的超平面时, 它在 $H_0$ 处与 $\check X$ 相切: 换言之, $H_0$ 在 $x_0$ 处与 $X$ 相切蕴含 $x_0$ 在 $H_0$ 处与 $\check X$ 相切. 由于 $\varphi$ 在 $(x_0,H_0)$ 处非分歧, 且 $\check X$ 在 $H_0=\varphi(x_0,H_0)$ 处光滑, 其微分是同构
\begin{equation}\tag{3.5.4}
d\varphi:
T_{\mathbb P(N),(x_0,H_0)}\xrightarrow{\sim}T_{\check X,H_0}.
\end{equation}
因此以下引理完成证明的第一部分.